\SourceLeaf{169}die Summe aller Annäherungen, so viele male man sie auch fortsetzte,
doch immer kleiner als $\frac{1}{1-\psi} \times ckf$ bleiben, und
diese Grenze könnte denn immer noch kleiner als der rechte Winkel
$dk\varphi$ sein. Du hast mein aufrichtiges Urtheil verlangt: ich
habe es gegeben, und ich wiederhole nochmals die Versicherung, dass
es mich innig freuen soll, wenn Du alle Schwierigkeiten überwindest.

.\, .\, .\, .\, .\, .\, .\, .\, .\, .\, .\, .\, .\, .\, .\, .\, .

\medskip
\begin{center}
\rule{0.3\textwidth}{0.4pt}
\end{center}
\medskip

\begin{center}
\textsc{Bemerkungen.}
\end{center}

\textsc{Gauss} und \textsc{Wolfgang Bolyai} aus Bolya in
Siebenbürgen haben vom Herbst 1796 bis Herbst 1798 zusammen in
Göttingen studirt und freundschaftlich verkehrt. Vor
\textsc{Bolyais} Heimreise sind sie noch einmal am 25. Mai 1799 in
Clausthal zusammengetroffen; der Brief vom 16. December 1799, von
dem hier ein Stück mitgetheilt wird, ist der erste, den \textsc{Gauss}
nach dem Abschiede an seinen Freund gerichtet hat.

\textsc{Bolyai} hat seine »Göttingische Theorie der Parallelen« erst im Jahre 1804 an \textsc{Gauss} geschickt, dessen
Kritik der Brief vom 25. November 1804 enthält.

Was \textsc{Gauss}’ Untersuchungen über die ersten Gründe der
Geometrie betrifft, so werden die vorstehenden Äusserungen ergänzt
durch die Notizen:

\begin{quote}
\emph{»Plani possibilitatem demonstravi. [1797] Jul. 28. Gotting.«}
\end{quote}
und
\begin{quote}
\emph{»In principiis Geometriae egregios progressus fecimus.
Br[unovici 1799] Sept.«}
\end{quote}
die sich in dem bereits S. 20 dieses Bandes erwähnten \textls{Tagebuche}
finden.

\begin{flushright}
\textsc{Stäckel.}
\end{flushright}

\medskip
\begin{center}
\rule{0.3\textwidth}{0.4pt}
\end{center}
